\begin{examplebox}
\textbf{\textcolor{CExample}{导入例：初中圆的水平弦，先不使用高中公式}}

\textbf{题目：}在圆 $x^2+y^2=25$ 中，直线 $y=3$ 截得弦 $AB$，求 $|AB|$。

\textbf{初中法（垂径定理 $+$ 勾股）：}
圆心为 $O(0,0)$，半径 $r=5$，圆心到弦的距离 $d=3$。
设弦中点为 $M$，则 $OM\perp AB$、$AM=MB$，故
\[
 AM=\sqrt{OA^2-OM^2}=\sqrt{25-9}=4,
 \qquad\boxed{|AB|=8.}
\]
\textbf{坐标法（发现两种方法相通）：}
联立 $y=3$ 得 $x^2=16$，于是
$A(-4,3),B(4,3)$。该直线水平，$k=0$，因此
\[
 |AB|=\sqrt{1+0^2}\,|{-4}-4|=8.
\]
\textbf{桥接：}同一段弦，同一个长度；两种不同的直角三角形都使用勾股定理。前一种构造“圆半径—圆心距—半弦”的直角三角形，
后一种直接用“水平坐标差—竖直坐标差—整弦”的直角三角形。

\medskip
\textbf{\textcolor{CExample}{过渡例：同一条圆斜弦，用初中法和高中法各算一次}}

\textbf{题目：}圆 $x^2+y^2=25$ 上取
$A(5,0)$、$B(-3,4)$，求弦 $AB$ 的长度，并解释两种算法为何一致。

\begin{center}
\begin{tikzpicture}[font=\small]
 \path[use as bounding box] (-3.3,-3.3) rectangle (4,3.5);
 \node at (.2,3.2) {同一条圆斜弦，两种直角三角形};
 \draw[blue!60!black,thick] (0,0) circle (2.5);
 \coordinate (A) at (2.5,0); \coordinate (B) at (-1.5,2);
 \coordinate (M) at (.5,1); \coordinate (H) at (-1.5,0);
 \draw[magenta!70!black,very thick] (A)--(B);
 \draw[teal!70!black,thick] (0,0)--(B);
 \draw[teal!70!black,dashed] (0,0)--(M);
 \draw[orange!80!black,dashed,thick] (A)--(H)--(B);
 \draw ($(M)+(-.16,.08)$)--++(-.08,-.16)--++(.16,-.08);
 \draw ($(H)+(.18,0)$)--++(0,.18)--++(-.18,0);
 \node[right] at (A) {$A(5,0)$}; \node[above left] at (B) {$B(-3,4)$};
 \node[below] at (0,0) {$O$}; \node[above] at (.5,1.45) {$M(1,2)$};
 \node[below] at (H) {$H$};
 \node at (0,-1) {$OM=d=\sqrt5$};
 \node[left] at (-1.5,.85) {$4$}; \node[below] at (.9,0) {$8$};
 \node at (.2,-3) {$OM\perp AB,\quad AH\perp BH,\quad |AB|=4\sqrt5$};
\end{tikzpicture}
\end{center}


\textbf{方法一（初中圆弦法）：}
经过两点的直线为 $x+2y-5=0$。
圆心 $O(0,0)$ 到该直线的距离
$\displaystyle d=\frac{|{-5}|}{\sqrt{1^2+2^2}}=\sqrt5$。
由垂径定理和勾股定理
\[
 |AB|=2\sqrt{r^2-d^2}
 =2\sqrt{25-5}=\boxed{4\sqrt5}.
\]

\textbf{方法二（高中坐标投影法）：}
该直线的斜率 $k=-\frac12$，两点横坐标差的绝对值为
$|5-(-3)|=8$，所以
\[
 |AB|=\sqrt{1+\left(-\tfrac12\right)^2}\,8
 =\frac{\sqrt5}{2}\times8=\boxed{4\sqrt5}.
\]

\textbf{迁移启发：}计算的是\emph{同一条}弦，两种公式自然相等；
从圆转向椭圆、双曲线、抛物线后，\textbf{保留第二种“投影差 $+$ 勾股”的通用方法}。

\medskip
\textbf{\textcolor{CExample}{例一（基础：圆的斜弦）}}

\textbf{题目：}直线 $y=x+1$ 与圆 $x^2+y^2=5$ 交于 $A,B$，求 $|AB|$。

\textbf{识别信号：}$k=1$，联立能得到横坐标二次方程，先看判别式。
\[
 x^2+(x+1)^2=5\ \Longrightarrow\  x^2+x-2=0.
\]
因此 $\alpha=1,\beta=1,\gamma=-2$，
$\Delta=1-4(1)(-2)=9>0$，
\[
 \boxed{|AB|=\sqrt{1+1^2}\,\frac{\sqrt9}{1}=3\sqrt2.}
\]
\textbf{另一条验证路线：}韦达给出
$x_1+x_2=-1,\ x_1x_2=-2$，
$|x_1-x_2|=\sqrt{(-1)^2-4(-2)}=3$。

\medskip
\textbf{\textcolor{CExample}{例二（抛物线：优先纵坐标消元）}}

\textbf{题目：}抛物线 $y^2=2x$ 与 $y=2x-1$ 交于 $A,B$，求 $|AB|$。

\textbf{识别信号：}曲线天然以 $y^2$ 表示，改写
$x=(y+1)/2=\frac12 y+\frac12$，直接消去 $x$。
\[
 y^2=y+1\ \Longrightarrow\ y^2-y-1=0,
 \qquad \Delta_y=1+4=5>0.
\]
取 $t=\frac12$，纵坐标差
$|y_1-y_2|=\sqrt5$，得
\[
 \boxed{|AB|=\sqrt{1+\left(\frac12\right)^2}\sqrt5
 =\frac{\sqrt5}{2}\sqrt5=\frac52.}
\]
\textbf{省时点：}不必分别求 $y_1,y_2$，也不必将 $x=(y+1)/2$ 再代回去求点坐标。

\medskip
\textbf{\textcolor{CExample}{例三（综合：椭圆平行弦最值）}}

\textbf{题目：}椭圆 $\frac{x^2}{2}+y^2=1$ 与斜率为 $1$ 的直线
$l:y=x+m$ 交于 $A,B$。求 $|AB|$ 的最大值及相应的 $m$。

\begin{center}
\begin{tikzpicture}[x=2cm,y=2cm,font=\small]
 \path[use as bounding box] (-2.7,-1.6) rectangle (2.7,1.9);
 \node at (0,1.65) {椭圆中的同方向弦：过中心时最长};
 \draw[->,gray] (-1.65,0)--(1.65,0) node[right] {$x$};
 \draw[->,gray] (0,-1.25)--(0,1.25) node[above] {$y$};
 \draw[blue!60!black,thick] (0,0) ellipse ({sqrt(2)} and 1);
 \foreach \m/\col/\sty in {0/magenta!70!black/solid,.8/teal!70!black/dashed,-.8/orange!80!black/dotted}{
   \pgfmathsetmacro{\lo}{(-2*\m-sqrt(6-2*\m*\m))/3}
   \pgfmathsetmacro{\hi}{(-2*\m+sqrt(6-2*\m*\m))/3}
   \draw[very thick,\col,\sty] (\lo,{\lo+\m})--(\hi,{\hi+\m});
 }
 \node[left,teal!70!black] at (-1.6,.9) {$m=0.8$};
 \draw[teal!70!black,thin] (-1.55,.8)--(-.8,0);
 \node[right,orange!80!black] at (1.6,-.9) {$m=-0.8$};
 \draw[orange!80!black,thin] (1.55,-.8)--(.8,0);
 \node[right,magenta!70!black] at (1.25,1.05) {$m=0$};
 \draw[magenta!70!black,thin] (1.2,1)--(.65,.65);
 \node at (0,-1.5) {$|AB|=\frac43\sqrt{3-m^2},\quad |m|<\sqrt3$};
\end{tikzpicture}
\end{center}


\textbf{第一步（联立并检查交点）：}
\[
 \frac{x^2}{2}+(x+m)^2=1
 \ \Longrightarrow\ 3x^2+4mx+2(m^2-1)=0.
\]
\[
 \Delta=(4m)^2-4\cdot3\cdot2(m^2-1)
 =24-8m^2.
\]
有两个不同交点当且仅当
$\Delta>0\iff-\sqrt3<m<\sqrt3$。

\textbf{第二步（不用求根，直接算弦长）：}
\[
 \begin{aligned}
 |AB|&=\sqrt2\,\frac{\sqrt{24-8m^2}}3\\
 &=\frac43\sqrt{3-m^2}.
 \end{aligned}
\]
\textbf{第三步（最值）：}由于 $m^2\ge0$，
在 $m=0$ 时取等号，且 $m=0$ 满足两交点条件，
\[
 \boxed{|AB|_{\max}=\frac{4\sqrt3}{3},\qquad m=0.}
\]
\textbf{几何复核：}中心在原点的椭圆中，同方向的平行弦以过中心的弦最长；$m=0$ 与结论相符。

\medskip
\textbf{\textcolor{CExample}{例四（特殊线：垂直于 $x$ 轴）}}

\textbf{题目：}双曲线 $\frac{x^2}{4}-y^2=1$ 被直线 $x=3$ 截得弦长多少？

\textbf{识别信号：}竖直线，斜率不存在；不要使用含 $k$ 的形式。
\[
 \frac94-y^2=1\ \Longrightarrow\ y^2=\frac54
 \ \Longrightarrow\ y_{1,2}=\pm\frac{\sqrt5}{2}.
\]
因此
\[
 \boxed{|AB|=|y_1-y_2|=\sqrt5.}
\]
\textbf{说明：}竖直线可看作 $x=0\cdot y+3$；统一纵坐标式系数为 $1$。

\tcbbreak
\textbf{\textcolor{CExample}{例五（高价值反例：退化成一次方程）}}

\textbf{题目：}双曲线 $x^2-y^2=1$ 与 $y=x+1$ 能否构成两个交点间的弦？

\begin{center}
\begin{tikzpicture}[font=\small]
 \path[use as bounding box] (-4,-3.4) rectangle (4,4.1);
 \node at (0,3.8) {与渐近线平行：只有一个有限交点};
 \draw[->,gray] (-3.4,0)--(3.4,0) node[right] {$x$};
 \draw[->,gray] (0,-2.6)--(0,3.1) node[above] {$y$};
 \draw[gray,dashed] (-2.8,-2.8)--(2.8,2.8) node[right] {$y=x$};
 \foreach \sgn in {-1,1} {
 \draw[blue!60!black,thick,domain=-2.5:2.5,samples=90] plot ({\sgn*sqrt(1+\x*\x)},\x);
 }
 \draw[magenta!70!black,thick] (-3.3,-2.3)--(2,3) node[above right] {$y=x+1$};
 \fill (-1,0) circle (2pt) node[above left] {$(-1,0)$};
 \node at (0,-3.1) {$x^2-y^2=1\quad\Longrightarrow\quad -2x-2=0$};
\end{tikzpicture}
\end{center}


联立得到
\[
 x^2-(x+1)^2=1
 \ \Longrightarrow\ -2x-2=0
 \ \Longrightarrow\ (x,y)=(-1,0).
\]
只存在\textbf{一个有限实交点}。原因是直线方向与渐近线 $y=x$ 平行，
联立后 $x^2$ 项抵消了。
\[
 \boxed{\text{没有两不同交点，不能使用二次方程弦长公式。}}
\]
\textbf{高考检查项：}先判断二次项系数是否非零，再使用 $\Delta>0$。
\end{examplebox}
